%\chapter*{Введение}
\ssr{Введение}

Целью данной лабораторной работы является:
\begin{itemize}
	\item изучить методы анализа структуры ориентированных и неориентированных графов;
	\item освоить поиск компонент сильной связности, мостов, точек сочленения и топологическую сортировку;
	\item реализовать и сравнить основные алгоритмы поиска кратчайших путей во взвешенных графах.
\end{itemize}

Задачи:
\begin{enumerate}
	\item научиться представлять граф списком смежности;
	\item научиться применять DFS как основу для структурного анализа;
	\item научиться использовать времена входа tin и значения low;
	\item реализовать алгоритм Дейкстры с приоритетной очередью, алгоритм Беллмана–Форда и алгоритм Флойда–Уоршелла;
	\item научиться определять условия применимости каждого метода.
\end{enumerate}

\ssr{Выполнение}

\ssection{Задание 1. Представление графа и базовый обход}

\begin{figure}[!htbp]
	\centering
	%\includegraphics[width=0.6\linewidth]{""}
	\caption{}
\end{figure}
\FloatBarrier

%\lstinputlisting[
%	\label{key}
%	style=python, 
%	caption={Список импортов и константы},
%	linerange={1-9}
%]{main.py}

\ssection{Задание 2. Компоненты сильной связности}

\ssection{Задание 3. Мосты и точки сочленения}

\ssection{Задание 4. Граф конденсации, топологическая сортировка и цикл}

\ssection{Задание 5. Алгоритм Дейкстры и восстановление маршрута}

\ssection{Задание 6. Беллман–Форд и отрицательный цикл}

\ssection{Задание 7. Флойд–Уоршелл и итоговое сравнение}

\ssr{Вывод}

В ходе лабораторной работы были получены:
\begin{itemize}
	\item знания, как представлять граф списком смежности;
	\item знания, как применять DFS как основу для структурного анализа;
	\item знания, как использовать времена входа tin и значения low;
	\item реализации алгоритма Дейкстры с приоритетной очередью, алгоритма Беллмана–Форда и алгоритма Флойда–Уоршелла;
	\item знания, как определять условия применимости каждого метода.
\end{itemize}

\ssr{Контрольные вопросы}
\begin{enumerate}
	\item \textit{Что такое список смежности и какова стоимость обхода графа при таком представлении?}\newline
	
	\item \textit{Чем компонента сильной связности отличается от обычной компоненты связности?}\newline
	
	\item \textit{Из каких основных этапов состоит алгоритм Косарайю?}\newline
	
	\item \textit{Какую роль играют tin и low при DFS?}\newline
	
	\item \textit{Как определяется мост через значения tin и low?}\newline
	
	\item \textit{Что называется точкой сочленения и чем отличается обработка корня DFS?}\newline
	
	\item \textit{Что такое DAG?}\newline
	
	\item \textit{Что означает топологический порядок и когда он существует?}\newline
	
	\item \textit{Как с помощью DFS обнаружить цикл в ориентированном графе?}\newline
	
	\item \textit{Что называется релаксацией ребра?}\newline
	
	\item \textit{При каких условиях корректен алгоритм Дейкстры?}\newline
	
	\item \textit{Зачем в Дейкстре используется приоритетная очередь?}\newline
	
	\item \textit{Как восстановить маршрут после поиска кратчайших расстояний?}\newline
	
	\item \textit{Почему отрицательные рёбра создают проблему для алгоритма Дейкстры?}\newline
	
	\item \textit{Как работает алгоритм Беллмана–Форда?}\newline
	
	\item \textit{Почему для кратчайшего пути без отрицательных циклов достаточно V-1 итераций релаксации?}\newline
	
	\item \textit{Как Беллман–Форд обнаруживает отрицательный цикл?}\newline
	
	\item \textit{Какую задачу решает алгоритм Флойда–Уоршелла?}\newline
	
	\item \textit{Какова временная сложность Дейкстры с двоичной кучей, Беллмана–Форда и Флойда–Уоршелла?}\newline
	
	\item \textit{В каких случаях следует выбирать каждый из трёх алгоритмов кратчайших путей?}\newline
	
\end{enumerate}